\chapter{Bad Breath (Halitosis)}

\section*{Definition}
Unpleasant odor of the breath, caused primarily by volatile sulfur compounds produced by bacterial metabolism in the oral cavity, and occasionally from systemic conditions.

\section*{When to See a Doctor}

\begin{itemize}
 \item See a dentist if bad breath persists for a few weeks despite careful self-care, or if you have toothache, bleeding or swollen gums, loose teeth, or denture problems.
 \item Seek medical assessment for bad breath with fever, facial swelling, difficulty swallowing or breathing, or other systemic symptoms.
\end{itemize}

\section*{How It Develops}
Most persistent halitosis originates in the mouth. Bacteria in tongue coating and periodontal pockets break down proteins and produce volatile sulfur compounds, including hydrogen sulfide and methyl mercaptan. Dry mouth reduces saliva's cleansing effect. Less commonly, the source is in the nose or throat, or a medical condition; gastrointestinal disease is an uncommon cause.

\section*{Causes}
\begin{itemize}
 \item Poor oral hygiene (tongue coating, food debris)
 \item Periodontal disease (gingivitis, periodontitis)
 \item Tonsillitis or tonsil stones (tonsilloliths)
 \item Dry mouth (xerostomia)
 \item Dental caries or abscesses
 \item Sinusitis or postnasal drip
 \item Gastroesophageal reflux disease (GERD)
 \item Strong foods (garlic, onions, spices)
 \item Smoking or tobacco use
 \item Systemic diseases such as poorly controlled diabetes, advanced liver disease, or advanced kidney disease (uncommon causes)
\end{itemize}

\section*{Related Symptoms}
\begin{itemize}
 \item Dry mouth
 \item Dental pain
 \item Postnasal drip
 \item Sore throat
 \item Acid reflux
\end{itemize}

\section*{Medical Evaluation}
\begin{itemize}
 \item Your doctor will examine the oral cavity, oropharynx, and nasal passages to identify dental caries, periodontal disease, tonsilloliths, or sinusitis.
 \item The clinician or dentist asks about oral hygiene, dry mouth, medicines, diet, tobacco, reflux, nasal symptoms, and how long the odor has been present, then examines the mouth, teeth, gums, tongue, throat, and nose as indicated.
 \item Breath-odor meters are not routinely needed and do not replace a clinical examination. Dental X-rays are ordered only when examination suggests hidden decay, infection, or bone loss.
 \item Dental assessment is usually the first referral for persistent halitosis; medical or ear, nose, and throat assessment is guided by the findings.
\end{itemize}

\section*{Treatment Options}
\begin{itemize}
 \item Professional dental cleaning and treatment of decay or gum disease, together with improved oral hygiene, usually address halitosis of oral origin.
 \item Periodontal treatment may include professional debridement and, when indicated, a clinician-directed antimicrobial mouthrinse; chlorhexidine can stain teeth and alter taste and should not be used long term without dental advice.
 \item Antibiotics, antifungals, or other medicines are used only when a specific infection or condition has been diagnosed; they are not routine treatment for bad breath.
 \item Treatment of an identified non-oral cause, such as significant dry mouth, nasal disease, reflux, or poorly controlled diabetes, is tailored to that cause.
\end{itemize}

\section*{Self-Care and Home Management}
\begin{itemize}
 \item Brush teeth twice daily with fluoride toothpaste and floss once daily to remove food debris and plaque between teeth.
 \item Clean the tongue dorsum daily using a tongue scraper or toothbrush to reduce the bacterial biofilm responsible for volatile sulfur compounds.
 \item Drink water regularly throughout the day to maintain saliva production; chew sugar-free gum or suck on sugar-free mints if dry mouth is a problem.
 \item Avoid tobacco products and limit alcohol. Strong-smelling foods can cause temporary odor but do not usually explain persistent halitosis.
 \item Keep dentures clean and remove them at night if you wear them. Replace a toothbrush when it is worn, and attend dental check-ups as recommended for your oral health.
\end{itemize}

\section*{References}
\begin{thebibliography}{99}
\bibitem{bad_breath:ref1} Hughes FJ, McNab R. ``Halitosis.'' \textit{BMJ}. 2010;340:c2191.
\bibitem{bad_breath:ref2} Torsten M, Gomez SM, et al. ``Oral malodour: an interdisciplinary approach.'' \textit{Br Dent J}. 2012;213(6):299-304.
\bibitem{bad_breath:ref3} Scully C, Greenman J. ``Halitosis (breath odor).'' \textit{Periodontol 2000}. 2008;48:66-75.
\bibitem{bad_breath:ref4} Newman MG, Takei HH, et al. \textit{Newman and Carranza\'s Clinical Periodontology}, 13th ed. Elsevier, 2018.
\bibitem{bad_breath:ref5} National Health Service. ``Bad breath.'' Reviewed June 2025. https://www.nhs.uk/symptoms/bad-breath/.
\bibitem{bad_breath:ref6} American Dental Association. ``Mouthrinse (Mouthwash).'' Accessed September 2026. https://www.ada.org/resources/ada-library/oral-health-topics/mouthrinse-mouthwash.
\end{thebibliography}
